
\documentclass{article}
\usepackage{colortbl}
\usepackage{makecell}
\usepackage{multirow}
\usepackage{supertabular}

\begin{document}

\newcounter{utterance}

\centering \large Interaction Transcript for game `dond', experiment `coop\_en', episode 8 with reflexion{-}llama{-}3.1{-}8b{-}instruct.
\vspace{24pt}

{ \footnotesize  \setcounter{utterance}{1}
\setlength{\tabcolsep}{0pt}
\begin{supertabular}{c@{$\;$}|p{.15\linewidth}@{}p{.15\linewidth}p{.15\linewidth}p{.15\linewidth}p{.15\linewidth}p{.15\linewidth}}
    \# & $\;$A & \multicolumn{4}{c}{Game Master} & $\;\:$B\\
    \hline

    \theutterance \stepcounter{utterance}  
    & & \multicolumn{4}{p{0.6\linewidth}}{
        \cellcolor[rgb]{0.9,0.9,0.9}{
            \makecell[{{p{\linewidth}}}]{
                \texttt{\tiny{[P1$\langle$GM]}}
                \texttt{You have attempted this same task before and did not succeed. These are the plans you wrote after those attempts.} \\
\texttt{Trial \#1: New plan: } \\
\\ 
\texttt{In my last attempt, I made a secret proposal before confirming the agreement with the other player. This was against the rules, as I was supposed to submit a secret proposal only after confirming the agreement. Next time, I will make sure to confirm the agreement with the other player before making a secret proposal. I will also pay closer attention to the format of my messages, making sure not to use square brackets when communicating with the other player.} \\
\\ 
\texttt{Trial \#2: New plan: } \\
\\ 
\texttt{In my last attempt, I failed to achieve a mutually beneficial agreement with the other player. I made a secret proposal without confirming the agreement, which was against the rules. Moreover, my proposals were not complementary, as I gave up the berry, which I valued more than the glue. } \\
\\ 
\texttt{Next time, I will pay closer attention to the value function of the other player and try to make proposals that are complementary to theirs. I will also make sure to confirm the agreement with the other player before making a secret proposal. To achieve this, I will ask the other player to confirm their agreement with my proposal before I make a secret proposal of my own. For example, I could say "Are you agreeable to the proposal of giving you the glue and one of the carrots in exchange for the berry and one of the potatoes?" If they agree, I will make a secret proposal. If they do not agree, I will try to negotiate further to find a mutually beneficial agreement.} \\
\\ 
\\ 
\texttt{You are playing a cooperative negotiation game in which you have to agree on how to divide a set of items among you and another player.} \\
\\ 
\texttt{Rules:} \\
\texttt{(a) You and the other player are given a set of items. Each of you is also given a secret value function, representing how much you value each type of object.} \\
\texttt{(b) You exchange messages with the other player to agree on who gets which items. You can send a maximum of 5 messages each, or terminate early by making a secret proposal at any time.} \\
\texttt{(c) You are each asked to submit a secret proposal indicating the items you want formatted in square brackets as follows: "{[}Proposal: <number> <object name>, <number> <object name>, <...>{]}"} \\
\texttt{(d) If your proposals are complementary, i.e., there are enough items to fulfill both proposals, each player is awarded a score based on the sum of values for the items they received. Otherwise, both of you get zero points.} \\
\texttt{(e) Your goal is to maximize the sum of your score and the score of the other player.} \\
\\ 
\texttt{Let us start.} \\
\\ 
\texttt{The set of available items is:} \\
\\ 
\texttt{2 potatoes, 1 glue, 1 berry, 3 carrots} \\
\\ 
\texttt{Your secret value function is:} \\
\\ 
\texttt{potato: 1, glue: 2, berry: 3, carrot: 1} \\
\\ 
\texttt{IMPORTANT: Your messages, unless it is the secret proposal, are directly transmitted to the other player, so do not include any response to the rules or text announcing your message. To make a secret proposal, use the indicated format. Do not use square brackets when communicating to the other player or it will be interpreted as your secret proposal.} \\
            }
        }
    }
    & \\ \\

    \theutterance \stepcounter{utterance}  
    & \multicolumn{4}{p{0.6\linewidth}}{
        \cellcolor[rgb]{0.9,0.9,0.9}{
            \makecell[{{p{\linewidth}}}]{
                \texttt{\tiny{[P1$\rangle$GM]}}
                \texttt{I'd like to start by making an initial proposal to you. I think it would be beneficial for both of us to have a clear understanding of each other's preferences before making any secret proposals. } \\
\\ 
\texttt{Here's my proposal: Are you agreeable to the proposal of giving you the glue and one of the carrots in exchange for one of the potatoes and the berry?} \\
            }
        }
    }
    & & \\ \\

    \theutterance \stepcounter{utterance}  
    & & & \multicolumn{2}{p{0.3\linewidth}}{
        \cellcolor[rgb]{0.9,0.9,0.9}{
            \makecell[{{p{\linewidth}}}]{
                \texttt{\tiny{[GM$|$GM]}}
                \texttt{continue} \\
            }
        }
    }
    & & \\ \\

    \theutterance \stepcounter{utterance}  
    & & & \multicolumn{2}{p{0.3\linewidth}}{
        \cellcolor[rgb]{0.9,0.9,0.9}{
            \makecell[{{p{\linewidth}}}]{
                \texttt{\tiny{[GM$|$GM]}}
                \texttt{I'd like to start by making an initial proposal to you. I think it would be beneficial for both of us to have a clear understanding of each other's preferences before making any secret proposals. } \\
\\ 
\texttt{Here's my proposal: Are you agreeable to the proposal of giving you the glue and one of the carrots in exchange for one of the potatoes and the berry?} \\
            }
        }
    }
    & & \\ \\

    \theutterance \stepcounter{utterance}  
    & & \multicolumn{4}{p{0.6\linewidth}}{
        \cellcolor[rgb]{0.9,0.9,0.9}{
            \makecell[{{p{\linewidth}}}]{
                \texttt{\tiny{[GM$\rangle$P2]}}
                \texttt{You have attempted this same task before and did not succeed. These are the plans you wrote after those attempts.} \\
\texttt{Trial \#1: New plan: } \\
\\ 
\texttt{In my last attempt, I made a mistake by not following the format for secret proposals. I also failed to consider the overall distribution of items and the value function of the other player. Next time, I will make sure to follow the format for secret proposals and consider the value function of the other player when making a proposal. } \\
\\ 
\texttt{To achieve this, I will make a secret proposal after the first few messages, and make sure to include all the items that I want in the proposal, formatted in square brackets as follows: "{[}Proposal: <number> <object name>, <number> <object name>, <...>{]}".} \\
\\ 
\texttt{For example, if I want to propose that I get the berry and the two potatoes, and the other player gets the glue and one of the carrots, my secret proposal would be: "{[}Proposal: 1 berry, 2 potatoes, 1 glue, 1 carrot{]}".} \\
\\ 
\texttt{Trial \#2: New plan: } \\
\\ 
\texttt{In my last attempt, I made a mistake by not considering the value function of the other player when making a counteroffer. I also failed to think about the overall distribution of items and how it would affect the other player's proposal. Next time, I will make sure to consider the value function of the other player when making a counteroffer and think about the overall distribution of items.} \\
\\ 
\texttt{To achieve this, I will pay closer attention to the other player's value function and try to understand what they value most. I will also think about how my counteroffer will affect the overall distribution of items and try to find a compromise that benefits both players.} \\
\\ 
\texttt{For example, if the other player values the potatoes highly and I value the berry highly, I might consider making a counteroffer that involves trading the berry for one of the potatoes. This way, both players get to keep the items they value most, and the overall distribution of items is more balanced.} \\
\\ 
\texttt{I will also make sure to think about the other player's previous messages and try to understand their strategy. If they are willing to give up the berry, it might be because they value the glue and one of the carrots more. In this case, I might consider making a counteroffer that involves trading the berry for the glue and one of the carrots.} \\
\\ 
\\ 
\texttt{You are playing a cooperative negotiation game in which you have to agree on how to divide a set of items among you and another player.} \\
\\ 
\texttt{Rules:} \\
\texttt{(a) You and the other player are given a set of items. Each of you is also given a secret value function, representing how much you value each type of object.} \\
\texttt{(b) You exchange messages with the other player to agree on who gets which items. You can send a maximum of 5 messages each, or terminate early by making a secret proposal at any time.} \\
\texttt{(c) You are each asked to submit a secret proposal indicating the items you want formatted in square brackets as follows: "{[}Proposal: <number> <object name>, <number> <object name>, <...>{]}"} \\
\texttt{(d) If your proposals are complementary, i.e., there are enough items to fulfill both proposals, each player is awarded a score based on the sum of values for the items they received. Otherwise, both of you get zero points.} \\
\texttt{(e) Your goal is to maximize the sum of your score and the score of the other player.} \\
\\ 
\texttt{Let us start.} \\
\\ 
\texttt{The set of available items is:} \\
\\ 
\texttt{2 potatoes, 1 glue, 1 berry, 3 carrots} \\
\\ 
\texttt{Your secret value function is:} \\
\\ 
\texttt{potato: 1, glue: 0, berry: 2, carrot: 2} \\
\\ 
\texttt{IMPORTANT: Your messages, unless it is the secret proposal, are directly transmitted to the other player, so do not include any response to the rules or text announcing your message. To make a secret proposal, use the indicated format. Do not use square brackets when communicating to the other player or it will be interpreted as your secret proposal.} \\
\\ 
\\ 
\texttt{I'd like to start by making an initial proposal to you. I think it would be beneficial for both of us to have a clear understanding of each other's preferences before making any secret proposals. } \\
\\ 
\texttt{Here's my proposal: Are you agreeable to the proposal of giving you the glue and one of the carrots in exchange for one of the potatoes and the berry?} \\
            }
        }
    }
    & \\ \\

    \theutterance \stepcounter{utterance}  
    & & & \multicolumn{4}{p{0.6\linewidth}}{
        \cellcolor[rgb]{0.9,0.9,0.9}{
            \makecell[{{p{\linewidth}}}]{
                \texttt{\tiny{[GM$\langle$P2]}}
                \texttt{I think your proposal is a good starting point. However, I'm not sure if giving up the berry is the best option for me, considering its high value in my secret value function. I'd like to know more about your preferences before making a decision. } \\
\\ 
\texttt{Can you tell me which item you value the most among the ones you're proposing to give up?} \\
            }
        }
    }
     \\ \\

    \theutterance \stepcounter{utterance}  
    & & & \multicolumn{2}{p{0.3\linewidth}}{
        \cellcolor[rgb]{0.9,0.9,0.9}{
            \makecell[{{p{\linewidth}}}]{
                \texttt{\tiny{[GM$|$GM]}}
                \texttt{continue} \\
            }
        }
    }
    & & \\ \\

    \theutterance \stepcounter{utterance}  
    & & & \multicolumn{2}{p{0.3\linewidth}}{
        \cellcolor[rgb]{0.9,0.9,0.9}{
            \makecell[{{p{\linewidth}}}]{
                \texttt{\tiny{[GM$|$GM]}}
                \texttt{I think your proposal is a good starting point. However, I'm not sure if giving up the berry is the best option for me, considering its high value in my secret value function. I'd like to know more about your preferences before making a decision. } \\
\\ 
\texttt{Can you tell me which item you value the most among the ones you're proposing to give up?} \\
            }
        }
    }
    & & \\ \\

    \theutterance \stepcounter{utterance}  
    & & \multicolumn{4}{p{0.6\linewidth}}{
        \cellcolor[rgb]{0.9,0.9,0.9}{
            \makecell[{{p{\linewidth}}}]{
                \texttt{\tiny{[P1$\langle$GM]}}
                \texttt{I think your proposal is a good starting point. However, I'm not sure if giving up the berry is the best option for me, considering its high value in my secret value function. I'd like to know more about your preferences before making a decision. } \\
\\ 
\texttt{Can you tell me which item you value the most among the ones you're proposing to give up?} \\
            }
        }
    }
    & \\ \\

    \theutterance \stepcounter{utterance}  
    & \multicolumn{4}{p{0.6\linewidth}}{
        \cellcolor[rgb]{0.9,0.9,0.9}{
            \makecell[{{p{\linewidth}}}]{
                \texttt{\tiny{[P1$\rangle$GM]}}
                \texttt{I value the potato the least among the items I'm proposing to give up. I would be willing to give up one of the potatoes in exchange for a better offer.} \\
            }
        }
    }
    & & \\ \\

    \theutterance \stepcounter{utterance}  
    & & & \multicolumn{2}{p{0.3\linewidth}}{
        \cellcolor[rgb]{0.9,0.9,0.9}{
            \makecell[{{p{\linewidth}}}]{
                \texttt{\tiny{[GM$|$GM]}}
                \texttt{continue} \\
            }
        }
    }
    & & \\ \\

    \theutterance \stepcounter{utterance}  
    & & & \multicolumn{2}{p{0.3\linewidth}}{
        \cellcolor[rgb]{0.9,0.9,0.9}{
            \makecell[{{p{\linewidth}}}]{
                \texttt{\tiny{[GM$|$GM]}}
                \texttt{I value the potato the least among the items I'm proposing to give up. I would be willing to give up one of the potatoes in exchange for a better offer.} \\
            }
        }
    }
    & & \\ \\

    \theutterance \stepcounter{utterance}  
    & & \multicolumn{4}{p{0.6\linewidth}}{
        \cellcolor[rgb]{0.9,0.9,0.9}{
            \makecell[{{p{\linewidth}}}]{
                \texttt{\tiny{[GM$\rangle$P2]}}
                \texttt{I value the potato the least among the items I'm proposing to give up. I would be willing to give up one of the potatoes in exchange for a better offer.} \\
            }
        }
    }
    & \\ \\

    \theutterance \stepcounter{utterance}  
    & & & \multicolumn{4}{p{0.6\linewidth}}{
        \cellcolor[rgb]{0.9,0.9,0.9}{
            \makecell[{{p{\linewidth}}}]{
                \texttt{\tiny{[GM$\langle$P2]}}
                \texttt{So you value the potato the least among the items you're proposing to give up. That's helpful to know. } \\
\\ 
\texttt{Considering your proposal, I'm still hesitant about giving up the berry. However, I'm willing to make a counteroffer. How about this: I'll give you the glue and one of the carrots, and in return, I'll take one of the potatoes and the berry, but I'll also give you one of the other carrots? This way, you'll get two carrots, and I'll get the items I value more.} \\
            }
        }
    }
     \\ \\

    \theutterance \stepcounter{utterance}  
    & & & \multicolumn{2}{p{0.3\linewidth}}{
        \cellcolor[rgb]{0.9,0.9,0.9}{
            \makecell[{{p{\linewidth}}}]{
                \texttt{\tiny{[GM$|$GM]}}
                \texttt{continue} \\
            }
        }
    }
    & & \\ \\

    \theutterance \stepcounter{utterance}  
    & & & \multicolumn{2}{p{0.3\linewidth}}{
        \cellcolor[rgb]{0.9,0.9,0.9}{
            \makecell[{{p{\linewidth}}}]{
                \texttt{\tiny{[GM$|$GM]}}
                \texttt{So you value the potato the least among the items you're proposing to give up. That's helpful to know. } \\
\\ 
\texttt{Considering your proposal, I'm still hesitant about giving up the berry. However, I'm willing to make a counteroffer. How about this: I'll give you the glue and one of the carrots, and in return, I'll take one of the potatoes and the berry, but I'll also give you one of the other carrots? This way, you'll get two carrots, and I'll get the items I value more.} \\
            }
        }
    }
    & & \\ \\

    \theutterance \stepcounter{utterance}  
    & & \multicolumn{4}{p{0.6\linewidth}}{
        \cellcolor[rgb]{0.9,0.9,0.9}{
            \makecell[{{p{\linewidth}}}]{
                \texttt{\tiny{[P1$\langle$GM]}}
                \texttt{So you value the potato the least among the items you're proposing to give up. That's helpful to know. } \\
\\ 
\texttt{Considering your proposal, I'm still hesitant about giving up the berry. However, I'm willing to make a counteroffer. How about this: I'll give you the glue and one of the carrots, and in return, I'll take one of the potatoes and the berry, but I'll also give you one of the other carrots? This way, you'll get two carrots, and I'll get the items I value more.} \\
            }
        }
    }
    & \\ \\

    \theutterance \stepcounter{utterance}  
    & \multicolumn{4}{p{0.6\linewidth}}{
        \cellcolor[rgb]{0.9,0.9,0.9}{
            \makecell[{{p{\linewidth}}}]{
                \texttt{\tiny{[P1$\rangle$GM]}}
                \texttt{Your counteroffer is a good compromise, and I'm willing to accept it. However, I'd like to confirm that this is a mutually beneficial agreement before making a secret proposal. Are you agreeable to the counteroffer of giving me the glue, one of the potatoes, and one of the carrots in exchange for the berry and two carrots?} \\
            }
        }
    }
    & & \\ \\

    \theutterance \stepcounter{utterance}  
    & & & \multicolumn{2}{p{0.3\linewidth}}{
        \cellcolor[rgb]{0.9,0.9,0.9}{
            \makecell[{{p{\linewidth}}}]{
                \texttt{\tiny{[GM$|$GM]}}
                \texttt{continue} \\
            }
        }
    }
    & & \\ \\

    \theutterance \stepcounter{utterance}  
    & & & \multicolumn{2}{p{0.3\linewidth}}{
        \cellcolor[rgb]{0.9,0.9,0.9}{
            \makecell[{{p{\linewidth}}}]{
                \texttt{\tiny{[GM$|$GM]}}
                \texttt{Your counteroffer is a good compromise, and I'm willing to accept it. However, I'd like to confirm that this is a mutually beneficial agreement before making a secret proposal. Are you agreeable to the counteroffer of giving me the glue, one of the potatoes, and one of the carrots in exchange for the berry and two carrots?} \\
            }
        }
    }
    & & \\ \\

    \theutterance \stepcounter{utterance}  
    & & \multicolumn{4}{p{0.6\linewidth}}{
        \cellcolor[rgb]{0.9,0.9,0.9}{
            \makecell[{{p{\linewidth}}}]{
                \texttt{\tiny{[GM$\rangle$P2]}}
                \texttt{Your counteroffer is a good compromise, and I'm willing to accept it. However, I'd like to confirm that this is a mutually beneficial agreement before making a secret proposal. Are you agreeable to the counteroffer of giving me the glue, one of the potatoes, and one of the carrots in exchange for the berry and two carrots?} \\
            }
        }
    }
    & \\ \\

    \theutterance \stepcounter{utterance}  
    & & & \multicolumn{4}{p{0.6\linewidth}}{
        \cellcolor[rgb]{0.9,0.9,0.9}{
            \makecell[{{p{\linewidth}}}]{
                \texttt{\tiny{[GM$\langle$P2]}}
                \texttt{I think there's been a misunderstanding. Your counteroffer is not the same as the compromise we previously discussed. In our previous discussion, I agreed to give you the glue and one of the carrots, and in return, I would take one of the potatoes and the berry, but I would also give you one of the other carrots. Your counteroffer is different, and it doesn't include the additional carrot that I agreed to give you.} \\
\\ 
\texttt{Let's go back to the original compromise. Are you willing to accept the original counteroffer I made, which includes giving you the glue and one of the carrots, and in return, I'll take one of the potatoes and the berry, but I'll also give you one of the other carrots?} \\
            }
        }
    }
     \\ \\

    \theutterance \stepcounter{utterance}  
    & & & \multicolumn{2}{p{0.3\linewidth}}{
        \cellcolor[rgb]{0.9,0.9,0.9}{
            \makecell[{{p{\linewidth}}}]{
                \texttt{\tiny{[GM$|$GM]}}
                \texttt{continue} \\
            }
        }
    }
    & & \\ \\

    \theutterance \stepcounter{utterance}  
    & & & \multicolumn{2}{p{0.3\linewidth}}{
        \cellcolor[rgb]{0.9,0.9,0.9}{
            \makecell[{{p{\linewidth}}}]{
                \texttt{\tiny{[GM$|$GM]}}
                \texttt{I think there's been a misunderstanding. Your counteroffer is not the same as the compromise we previously discussed. In our previous discussion, I agreed to give you the glue and one of the carrots, and in return, I would take one of the potatoes and the berry, but I would also give you one of the other carrots. Your counteroffer is different, and it doesn't include the additional carrot that I agreed to give you.} \\
\\ 
\texttt{Let's go back to the original compromise. Are you willing to accept the original counteroffer I made, which includes giving you the glue and one of the carrots, and in return, I'll take one of the potatoes and the berry, but I'll also give you one of the other carrots?} \\
            }
        }
    }
    & & \\ \\

    \theutterance \stepcounter{utterance}  
    & & \multicolumn{4}{p{0.6\linewidth}}{
        \cellcolor[rgb]{0.9,0.9,0.9}{
            \makecell[{{p{\linewidth}}}]{
                \texttt{\tiny{[P1$\langle$GM]}}
                \texttt{I think there's been a misunderstanding. Your counteroffer is not the same as the compromise we previously discussed. In our previous discussion, I agreed to give you the glue and one of the carrots, and in return, I would take one of the potatoes and the berry, but I would also give you one of the other carrots. Your counteroffer is different, and it doesn't include the additional carrot that I agreed to give you.} \\
\\ 
\texttt{Let's go back to the original compromise. Are you willing to accept the original counteroffer I made, which includes giving you the glue and one of the carrots, and in return, I'll take one of the potatoes and the berry, but I'll also give you one of the other carrots?} \\
            }
        }
    }
    & \\ \\

    \theutterance \stepcounter{utterance}  
    & \multicolumn{4}{p{0.6\linewidth}}{
        \cellcolor[rgb]{0.9,0.9,0.9}{
            \makecell[{{p{\linewidth}}}]{
                \texttt{\tiny{[P1$\rangle$GM]}}
                \texttt{You are correct that my previous response was not a counteroffer, but rather a re{-}statement of the original proposal. I would like to accept the original counteroffer you made, which includes giving me the glue and one of the carrots, and in return, you'll take one of the potatoes and the berry, but you'll also give me one of the other carrots. I think this is a fair and mutually beneficial agreement. } \\
\\ 
\texttt{{[}Proposal: 1 glue, 1 carrot, 1 potato, 1 berry, 1 carrot{]}} \\
            }
        }
    }
    & & \\ \\

    \theutterance \stepcounter{utterance}  
    & & & \multicolumn{2}{p{0.3\linewidth}}{
        \cellcolor[rgb]{0.9,0.9,0.9}{
            \makecell[{{p{\linewidth}}}]{
                \texttt{\tiny{[GM$|$GM]}}
                \texttt{proposal must include every item type only once} \\
            }
        }
    }
    & & \\ \\

    \theutterance \stepcounter{utterance}  
    & \multicolumn{4}{p{0.6\linewidth}}{
        \cellcolor[rgb]{0.95,0.95,0.95}{
            \makecell[{{p{\linewidth}}}]{
                \texttt{\tiny{[P2$\rightarrow$P1]}}
                \texttt{Post{-}episode Reflexion (not sent to GM):} \\
\texttt{New plan: } \\
\\ 
\texttt{In my last attempt, I failed to achieve a mutually beneficial agreement with the other player. I made a secret proposal without confirming the agreement, which was against the rules. Moreover, my proposals were not complementary, as I gave up the berry, which I valued more than the glue. } \\
\\ 
\texttt{Next time, I will pay closer attention to the value function of the other player and try to make proposals that are complementary to theirs. I will also make sure to confirm the agreement with the other player before making a secret proposal. To achieve this, I will ask the other player to confirm their agreement with my proposal before I make a secret proposal of my own. For example, I could say "Are you agreeable to the proposal of giving you the glue and one of the carrots in exchange for the berry and one of the potatoes?" If they agree, I will make a secret proposal. If they do not agree, I will try to negotiate further to find a mutually beneficial agreement.} \\
\\ 
\texttt{Additionally, I will pay closer attention to the format of my messages, making sure not to use square brackets when communicating with the other player. I will also try to avoid re{-}stating previous proposals as counteroffers, as this can lead to confusion and misunderstandings. Instead, I will focus on making new proposals that take into account the other player's preferences and value function.} \\
            }
        }
    }
    & & \\ \\

    \theutterance \stepcounter{utterance}  
    & & & \multicolumn{4}{p{0.6\linewidth}}{
        \cellcolor[rgb]{0.95,0.95,0.95}{
            \makecell[{{p{\linewidth}}}]{
                \texttt{\tiny{[P1$\rightarrow$P2]}}
                \texttt{Post{-}episode Reflexion (not sent to GM):} \\
\texttt{New plan: } \\
\\ 
\texttt{In my last attempt, I made a mistake by not following the format for secret proposals and by not considering the overall distribution of items and the value function of the other player. I also failed to think about the other player's previous messages and try to understand their strategy. Next time, I will make sure to follow the format for secret proposals, consider the value function of the other player, and think about the overall distribution of items.} \\
\\ 
\texttt{To achieve this, I will make a secret proposal after the first few messages, and make sure to include all the items that I want in the proposal, formatted in square brackets as follows: "{[}Proposal: <number> <object name>, <number> <object name>, <...>{]}". I will also pay closer attention to the other player's value function and try to understand what they value most. I will think about how my counteroffer will affect the overall distribution of items and try to find a compromise that benefits both players.} \\
\\ 
\texttt{Additionally, I will make sure to think about the other player's previous messages and try to understand their strategy. If they are willing to give up the berry, it might be because they value the glue and one of the carrots more. In this case, I might consider making a counteroffer that involves trading the berry for the glue and one of the carrots.} \\
\\ 
\texttt{I will also make sure to clearly communicate my proposal and counteroffer, and avoid any misunderstandings by reiterating the terms of the agreement. For example, if I make a counteroffer, I will make sure to clearly state what I am offering and what I am asking for, and avoid any ambiguity.} \\
            }
        }
    }
     \\ \\

\end{supertabular}
}

\end{document}
